% ch08-stale — the stale-covariate problem: formalization, damage, remedy ladder

\section{Why staleness owns a chapter}\label{ch08:sec:problem}

Defect (C5) of \cref{ch:data} is easy to state and easy to underestimate.
The firm covariates that every model in this document consumes ---
financials, headcount, capitalization --- are not observed as processes.
PitchBook records them when, and only when, the firm itself transacts:
each deal row carries the covariates as reported around that deal, and
between deals the pipeline carries the last recorded value forward
(\cref{eq:ch02:locf}, implemented in \code{synthetic/loaders.py} and to be
mirrored by every real-data loader per Protocol~\ref{prot:ch02:pit}). Two facts
make this more than a measurement-error nuisance. First, the error is not
noise but \emph{drift}: the gap between the last recorded value and the
current truth grows with the time since the firm's last deal. Second ---
and this is the point the whole chapter turns on --- the observation times
are event times \emph{of the process being modeled}. The firms that have
gone longest without a deal have the stalest covariates, and they are
exactly the low-intensity firms. Staleness is informative: the age of a
firm's covariates is itself evidence about the firm.

The rehearsal campaigns justify the chapter's rank. Design decision D2.1
of \cref{ch:data} records that the largest errors in the synthetic
rehearsal came not from wrong model families but from unstated exposure to
(C3) and (C5); the covariate-noise campaign (exp22, \cref{app:results})
measured the mechanics --- structure erodes before prediction does --- and
the standing verdict that \emph{covariates dominate, excitation is decimal
points} means the covariate channel is where the predictive mass lives.
A defect that silently corrupts the dominant channel, differentially
across the risk set, is one of this document's two reasons to exist.

The chapter proceeds in the document's standard order: formalize
(\cref{ch08:sec:formal}), prove the damage in a tractable case
(\cref{ch08:sec:consequences}), then commit to a remedy ladder --- R1,
staleness-aware observed-data modeling, committed and always on
(\cref{ch08:sec:r1}); R2, supervised drift modeling with regression
calibration (\cref{ch08:sec:r2}); R3, an explicit covariate state space
(\cref{ch08:sec:r3}); R4, joint modeling, adopted only if a defined
trigger fires (\cref{ch08:sec:r4}) --- and close with the bucket-model
face of the same mathematics, the rehearsal evidence, and the defect
declaration (\cref{ch08:sec:buckets,ch08:sec:declaration}).

\section{Formalization}\label{ch08:sec:formal}

Throughout, $i$ ranges over tracked firms, $N_i$ is firm $i$'s deal
counting process, $\Ft$ is the full filtration of
\cref{def:ch02:market} \emph{enlarged with the latent covariate paths},
and $\mathcal{G}_t \subset \Ft$ is the \emph{observed} filtration: the
internal history of the marked point process, the published macro paths,
and the covariate values recorded at past deals --- everything a
production pipeline can know at $t$, nothing it cannot.

\begin{assumption}[Latent covariate dynamics]\label{asm:ch08:latent}
Each tracked firm carries a latent c\`adl\`ag covariate process
$\xcov_i(t) \in \R^{q}$, partitioned into \emph{mechanical} components ---
deterministic functionals of the firm's own observed event history and
marks (raised-to-date, deal count, time-stamped last-round descriptors)
--- and \emph{drifting} components (revenue, EBITDA, headcount, latent
valuation). As a working model, the drifting components follow a
mean-reverting Ornstein--Uhlenbeck diffusion \citep{uhlenbeck1930} with
sector- and macro-driven mean,
\begin{equation}\label{eq:ch08:ou}
\dd \xcov_i(t)
 \;=\;
 -\Theta\,\big(\xcov_i(t) - m_i(t)\big)\dd t
 \;+\; \Sigma_W^{1/2} \dd W_i(t),
\qquad
 m_i(t) \;=\; m_{s(i)} + \Gamma\T \xcov^{\mathrm{mac}}(t),
\end{equation}
with $\Theta = \diag(\vartheta_1,\dots,\vartheta_q)$, $\vartheta_k > 0$
(the mean-reversion rates $\vartheta$ are unrelated to the excitation
decay $\decay$ of \cref{ch:hawkes}), $\Sigma_W$ diagonal, and $W_i$
idiosyncratic across firms given sector and macro. Where stationarity is
invoked, coordinate $k$ has stationary variance $\sigma_{X,k}^2 =
\Sigma_{W,kk}/(2\vartheta_k)$. Only the qualitative content --- mean
reversion toward an observable anchor, conditional variance growing in
elapsed time --- is load-bearing; the Gaussian linear form is what makes
the propositions below exact.
\end{assumption}

Read concretely, \cref{eq:ch08:ou} says: a firm's revenue and headcount
track an anchor set by its sector's trajectory and the macro state ---
rate cuts and sector booms lift the anchor, contractions lower it ---
while $\vartheta_k$ sets the memory of coordinate $k$, with covariate
half-life $\log 2 / \vartheta_k$. Nothing in the data pins these
half-lives a priori --- headcount is plausibly slow, revenue faster,
reported EBITDA noisier still --- and estimating them is exactly what the
drift dataset of \cref{ch08:sec:r2} is for. The OU form is chosen
because it is the tractable member of its qualitative class: it is the
Gaussian Markov process with mean reversion, it closes the filtering
recursions of \cref{ch08:sec:r3} in closed form, and it makes every
distortion formula in \cref{ch08:sec:consequences} exact rather than
first-order. When the working model fails --- heavy-tailed revenue jumps,
regime-switching growth --- the exact constants below degrade to
first-order approximations, but the monotone conclusions (attenuation
increasing in age, variance increasing in age, survival tilt downward)
survive, because they use only the qualitative content.

\begin{assumption}[LOCF observation scheme]\label{asm:ch08:locf}
The data delivers, for each firm and each $t$ past its entry $\entry_i$,
\begin{equation}\label{eq:ch08:locf}
\tilde\xcov_i(t) \;=\; \xcov_i\big(\tau_i(t)\big),
\qquad
\tau_i(t) \;=\; \sup\{t_m < t : i_m = i\},
\qquad
\staleage_i(t) \;=\; t - \tau_i(t),
\end{equation}
with the strict inequality of Protocol~\ref{prot:ch02:pit}(b): the value used at
$t$ comes from deals \emph{strictly before} $t$, so that
$\tilde\xcov_i$ and the staleness age $\staleage_i$ are left-continuous
in $t$. At its own deal times the firm's drifting covariates are observed
(as reported in the deal row); \cref{ch08:sec:r3} adds reporting noise.
For mechanical components LOCF is \emph{exact} --- they are step
functions that only move at own deals --- so (C5) binds on the drifting
components alone. This split is the basis of the column classification in
\cref{app:columns} and of decision D8.3 below.
\end{assumption}

\begin{proposition}[The observation scheme is internal and informative]
\label{prop:ch08:internal}
Under Assumptions~\ref{asm:ch08:latent} and~\ref{asm:ch08:locf}:
\begin{enumerate}[label=(\roman*)]
\item Each observation time of firm $i$'s covariates is a jump time of
$N_i$, hence a stopping time of the modeled filtration $\{\Ft\}$; the
processes $\tau_i(t)$, $\staleage_i(t)$, and $\tilde\xcov_i(t)$ are
left-continuous and adapted to $\{\mathcal{G}_t\}$, hence
$\mathcal{G}$-predictable --- they are legitimate intensity covariates.
\item The staleness age is a deterministic functional of the event
history: $\{\staleage_i(t) > u\} = \{N_i([t-u, t)) = 0\}$; it is the
backward recurrence time of the modeled process.
\item The scheme is \emph{informative}: whenever the true intensity
$\lam_i$ depends on $\xcov_i$, the conditional law of $\xcov_i(t)$ given
$(\tilde\xcov_i(t), \staleage_i(t))$ differs from the law obtained from
the free dynamics \cref{eq:ch08:ou} alone, because
$\Prob\big(\staleage_i(t) > u \st \F^{X}_t \vee \F_{\tau}\big) =
\exp\big(-\int_{t-u}^{t} \atrisk_i(v)\,\lam_i(v)\dd v\big)$ depends on
the latent covariate path over the gap.
\end{enumerate}
\end{proposition}

\begin{proof}
(i) Jump times of an adapted counting process are stopping times; between
jumps $\tau_i$ is frozen and $\staleage_i$ grows deterministically, and
the strict-inequality convention in \cref{eq:ch08:locf} makes all three
processes left-continuous adapted, hence predictable. (ii) is
\cref{eq:ch08:locf} read backwards. (iii) The displayed survival
probability is the standard conditional-inhomogeneous-Poisson formula
given the intensity path; if $\lam_i$ is $\xcov_i$-measurable and
nonconstant in it, the event $\{\staleage_i(t) > u\}$ has different
probability under different covariate paths, so by Bayes the conditional
law of the path given the gap is tilted (\cref{lem:ch08:tilt} computes
the tilt).
\end{proof}

In the Kalbfleisch--Prentice taxonomy of time-dependent covariates
\citep[Ch.~6]{kalbfleisch2002}, the macro series are \emph{external}:
they evolve regardless of the deal process, and survival-type
probability statements conditioning on their future paths are legitimate.
The pair $(\tilde\xcov_i, \staleage_i)$ is \emph{internal}: it is a
functional of the firm's own event history and marks. The consequence,
inherited verbatim from the internal-covariate discussion of
\cref{ch:hawkes}, is that intensities conditioned on
$(\tilde\xcov_i,\staleage_i)$ retain their interpretation, while plug-in
survival functionals (``probability of no deal in two years at today's
covariates'') do not --- the covariate would not stay at today's value.
One-step prediction is clean; multi-step statements must simulate the
system jointly, covariate observation process included.

How much information does the age itself carry? The rehearsal gives a
sharp, inherited answer (\cref{app:results}): in the model-ladder
campaign a recency-only ranker --- score firms by how recently they
raised --- performed \emph{worse than the popularity baseline}, because
in a cooldown world recent activity is anti-signal at short range; yet
the committed choice models carry an explicit gap covariate precisely
because the gap is informative --- with the measured caveat that its
fitted functional form is DGP-specific (on the misspecified regime the
gap weight tracked the generator's 30-day inhibition, not the cooldown
it was designed around). The age is neither ignorable nor
monotone in effect: it interacts with everything else, which is why
\cref{ch08:sec:r1} insists on interactions and why
\cref{ch08:sec:buckets} argues trees are the natural estimator of that
surface. The lesson to carry forward: $\staleage$ is simultaneously a
nuisance index (how corrupted is $\tilde\xcov$) and a signal (how
long has the market declined to fund this firm), and a correct treatment
must let the data express both roles at once.

\section{What staleness does to a log-linear intensity}
\label{ch08:sec:consequences}

Fix one firm, drop the subscript, and let the true intensity be
log-linear in the current latent covariates,
\begin{equation}\label{eq:ch08:truelam}
\lam(t) \;=\; \exp\!\big(\bcoef\T \xcov(t)\big),
\qquad
\xcov(t) \;=\; \tilde\xcov(t) + \eps(t),
\end{equation}
where $\eps(t) = \xcov(t) - \xcov(\tau)$ is the drift accumulated since
the last observation. (A baseline and external covariates can be carried
along untouched; they are omitted for legibility.) We fit with
$\tilde\xcov$ in place of $\xcov$. Everything in this section is a
statement about what that substitution does.

\subsection{Attenuation in the Gaussian-drift case}
\label{ch08:sec:attenuation}

Two benchmark drift laws bracket the mechanism. The first switches off
mean reversion and the survival information; it isolates the pure
variance effect.

\begin{proposition}[Independent Gaussian drift: multiplicative
distortion]\label{prop:ch08:berkson}
Suppose that, given $(\tilde\xcov, \staleage) = (\tilde x, a)$ and under
the free drift law (survival information ignored), $\eps(t) \sim
\mathcal{N}(0, \Sigma(a))$ independent of $\tilde x$ --- the
$\Theta \to 0$ (Brownian) limit of \cref{eq:ch08:ou}, where $\Sigma(a) =
\Sigma_W\, a$. Then the intensity implied for the observed data is
\begin{equation}\label{eq:ch08:berkson}
\tilde\lam(t)
 \;=\;
 \E\big[\exp(\bcoef\T\xcov(t)) \st \tilde x, a\big]
 \;=\;
 \exp\!\Big(\bcoef\T\tilde x \;+\; \tfrac12\,\bcoef\T\Sigma(a)\bcoef\Big):
\end{equation}
a staleness-dependent multiplicative distortion
$\exp\{\tfrac12\bcoef\T\Sigma(a)\bcoef\} \ge 1$, increasing in
$\staleage$, with the coefficient on $\tilde x$ unbiased at fixed $a$.
\end{proposition}

\begin{proof}
$\E[\exp(\bcoef\T\eps) \st \tilde x, a] =
\exp(\tfrac12\bcoef\T\Sigma(a)\bcoef)$ is the Gaussian moment generating
function; independence lets it factor out of
$\exp(\bcoef\T\tilde x)$.
\end{proof}

\begin{remark}
The structure ``truth $=$ observed $+$ independent error'' is
Berkson-type, and \cref{eq:ch08:berkson} shows its signature: no slope
attenuation \emph{conditional on} $\staleage$, but an omitted
$\staleage$-dependent term. Convexity of $\exp$ makes uncertain firms
look \emph{more} intense on average --- the variance term rewards
staleness --- which is one half of the ranking distortion of
\cref{prop:ch08:ranking}. The classical-error benchmark, $\tilde x =
\xcov + U$ with $U \perp \xcov$, gives the textbook attenuation instead:
$\E[\xcov \st \tilde x] = \mu + \Sigma_X(\Sigma_X + \Sigma_U)^{-1}(\tilde
x - \mu)$, so the induced log-intensity is linear in $\tilde x$ with
slope $\Sigma_X(\Sigma_X+\Sigma_U)^{-1}\bcoef$ --- in the scalar case
$\bcoef\,\sigma_X^2/(\sigma_X^2+\sigma_U^2)$
\citep{prentice1982,carroll2006}. The exp22 rehearsal campaign realizes
exactly this benchmark (\cref{ch08:sec:buckets}).
\end{remark}

The second benchmark restores mean reversion; it produces genuine,
staleness-graded attenuation and, as a corollary, the functional form of
remedy R1.

\begin{proposition}[Stationary OU drift: staleness-graded attenuation]
\label{prop:ch08:attenuation}
Let the drifting covariates follow \cref{eq:ch08:ou} with constant mean
$m$, diagonal $\Theta$, stationary initial condition, and survival
information ignored. Then, given $(\tilde x, a)$,
\begin{equation}\label{eq:ch08:oucond}
\xcov(t) \st \tilde x, a
 \;\sim\;
 \mathcal{N}\!\Big(m + e^{-\Theta a}(\tilde x - m),\; V(a)\Big),
\qquad
V(a)_{kk} = \sigma_{X,k}^2\big(1 - e^{-2\vartheta_k a}\big),
\end{equation}
and the observed-data intensity is
\begin{equation}\label{eq:ch08:attenuation}
\tilde\lam(t)
 \;=\;
 \exp\!\Big(
 \underbrace{\big(e^{-\Theta a}\bcoef\big)\T \tilde x}_{\text{attenuated
 signal}}
 \;+\;
 \underbrace{\bcoef\T\big(I - e^{-\Theta a}\big)m
 + \tfrac12\,\bcoef\T V(a)\bcoef}_{\;=\;g_0(a)\text{, staleness offset}}
 \Big).
\end{equation}
At fixed staleness the coefficient on $\tilde x_k$ is
$e^{-\vartheta_k a}\bcoef_k$: an attenuation factor $e^{-\vartheta_k a}
\in (0,1)$, decreasing in the age, equal to the regression-calibration
factor $\Cov(\xcov_k(t), \tilde x_k)/\Var(\tilde x_k)$.
\end{proposition}

\begin{proof}
Variation of constants on \cref{eq:ch08:ou} gives $\xcov(t) = m +
e^{-\Theta a}(\xcov(\tau) - m) + \int_\tau^t e^{-\Theta(t-u)}
\Sigma_W^{1/2}\dd W(u)$; the stochastic integral is Gaussian,
independent of $\xcov(\tau)$, with the stated variance (compute
$\int_0^a e^{-2\vartheta_k u}\Sigma_{W,kk}\dd u$ and use
$\sigma^2_{X,k} = \Sigma_{W,kk}/2\vartheta_k$), which proves
\cref{eq:ch08:oucond}. Then \cref{eq:ch08:attenuation} is the Gaussian
mgf as in \cref{prop:ch08:berkson}. For the last claim, stationarity
makes $(\xcov(\tau),\xcov(t))$ jointly Gaussian with
$\Var(\tilde x_k) = \sigma_{X,k}^2$ and $\Cov(\xcov_k(t),\tilde x_k) =
\sigma_{X,k}^2 e^{-\vartheta_k a}$, so the ratio is $e^{-\vartheta_k a}$;
equivalently, time reversibility of the stationary OU lets one write
$\tilde x_k = m_k + e^{-\vartheta_k a}(\xcov_k(t) - m_k) + U_k$ with
$U_k \perp \xcov_k(t)$ --- the classical shrink-plus-noise surrogate
structure of \citet{prentice1982}, with $e^{-\vartheta_k a}$ playing the
role of $\sigma_X^2/(\sigma_X^2 + \sigma_U^2)$.
\end{proof}

\begin{corollary}[Pooled fit: risk-set-weighted attenuation plus an
informativeness term]\label{cor:ch08:pooled}
Take the scalar case with centered covariate ($m = 0$), and fit the
misspecified model $\lam = \exp(c + b\,\tilde x)$ across the risk set
while the truth is \cref{eq:ch08:attenuation}. In the small-signal
(first-order) regime, the quasi-likelihood fit is the exposure-weighted
least-squares projection of the true log-intensity onto $(1, \tilde x)$,
so the estimand is
\begin{equation}\label{eq:ch08:pooled}
b
 \;=\;
 \frac{\Cov_w\!\big(\tilde x,\; e^{-\vartheta \staleage}\,\tilde x\big)}
      {\Var_w(\tilde x)}\;\bcoef
 \;+\;
 \frac{\Cov_w\!\big(\tilde x,\; g_0(\staleage)\big)}{\Var_w(\tilde x)},
\end{equation}
with weights $w$ proportional to exposure. If $\staleage \perp \tilde x$
across the risk set, $b = \E_w[e^{-\vartheta \staleage}]\,\bcoef$: pure
attenuation by the average factor. Informativeness breaks the
independence with a definite sign: if the intensity is increasing in
$\bcoef\T\xcov$ with $\bcoef > 0$, then $\staleage$ is stochastically
decreasing in $\tilde x$ (better last-seen firms raise sooner), so for
the increasing offset $g_0$ the second term is $\le 0$ --- the omitted
staleness structure biases the fitted slope \emph{further} toward zero,
beyond the average attenuation.
\end{corollary}

\begin{proof}
Linearizing the log-likelihood \cref{eq:ch02:jacod} in the small-signal
regime reduces the score equations for $(c,b)$ to the normal equations of
exposure-weighted least squares of the true log-intensity
$e^{-\vartheta a}\bcoef \tilde x + g_0(a)$ on $(1,\tilde x)$;
\cref{eq:ch08:pooled} is the resulting projection coefficient. Under
independence the first ratio is $\E_w[e^{-\vartheta a}]$ and the second
covariance vanishes. For the sign claim: $\staleage$ stochastically
decreasing in $\tilde x$ means $\E[g_0(\staleage)\st \tilde x]$ is
nonincreasing in $\tilde x$ for any increasing $g_0$, hence
$\Cov_w(\tilde x, g_0(\staleage)) \le 0$.
\end{proof}

This corollary is the precise version of the chapter's slogan: staleness
does not merely add noise, it attenuates the fitted covariate effects by
a factor tied to the age distribution of the risk set, and the
informativeness of the ages adds a signed omitted-variable term on top.
Both effects disappear from a model that conditions on the age --- which
is remedy R1.

\subsection{The informativeness twist: distortion is differential across
the risk set}\label{ch08:sec:twist}

\Cref{prop:ch08:berkson,prop:ch08:attenuation} deliberately ignored one
piece of information: the firm has \emph{not} raised during
$(\tau, t]$. By \cref{prop:ch08:internal}(iii) that silence is evidence.
The following lemma pins its direction; to stay tractable it summarizes
the intensity over the gap through the endpoint signal, which is
adequate for the slowly-varying covariates at issue.

\begin{lemma}[Survival tilt]\label{lem:ch08:tilt}
Let $S = \bcoef\T\xcov(t)$ have conditional density $p(s)$ given
$(\tilde x, a)$ under the free drift law, and suppose the integrated
intensity over the gap is $\int_\tau^t \lam(u)\dd u = a\,\Lambda(S)$
with $\Lambda$ strictly increasing (e.g.\ covariates approximately
constant over the gap, $\Lambda(s) = e^{s}$). Then the conditional
density given additionally $\{$no event in $(\tau, t]\}$ is
\[
p_{\mathrm{surv}}(s) \;\propto\; p(s)\, e^{-a\Lambda(s)},
\]
which is dominated by $p$ in the likelihood-ratio order; consequently,
for every increasing $h$ with finite means,
$\E[h(S) \st \tilde x, a, \text{no event}] \le \E[h(S) \st \tilde x, a]$,
strictly if $a > 0$, $\bcoef \ne 0$ and $S$ is nondegenerate.
\end{lemma}

\begin{proof}
Bayes gives the displayed density; the ratio
$p_{\mathrm{surv}}/p \propto e^{-a\Lambda(s)}$ is strictly decreasing in
$s$, which is likelihood-ratio dominance, and likelihood-ratio dominance
implies the expectation inequality for increasing $h$ (Chebyshev's
correlation inequality applied to the decreasing weight
$e^{-a\Lambda(s)}$ and increasing $h$). When the monotone-summary
assumption fails --- strongly oscillating paths --- the pointwise
ordering can fail locally, but the mechanism survives: the tilt
$\exp(-\int\lam)$ always downweights high-intensity \emph{paths}.
\end{proof}

\begin{proposition}[Sign of the ranking bias under naive LOCF plug-in]
\label{prop:ch08:ranking}
Rank firms at time $t$ by the naive plug-in score $\tilde s_i =
\bcoef\T\tilde\xcov_i(t)$, while the correct observed-data score (the
log of the projected intensity, \cref{prop:ch08:projection}) is $s_i^* =
\log\E[\exp(\bcoef\T\xcov_i(t)) \st \mathcal{G}_{t^-}]$. Under
Assumption~\ref{asm:ch08:latent} with stationary drift and \cref{lem:ch08:tilt}:
\begin{enumerate}[label=(\roman*)]
\item \emph{Conditional on the truth}: the score error is $\tilde s_i -
\bcoef\T\xcov_i(t) = -\bcoef\T\eps_i(t)$; a firm whose covariates
deteriorated in the $\bcoef$ direction since its last deal
($\bcoef\T\eps_i < 0$) is over-ranked, one that improved is
under-ranked, with error magnitude growing in the drift, hence
stochastically in $\staleage_i$.
\item \emph{Conditional on the data}: by
\cref{prop:ch08:attenuation} and \cref{lem:ch08:tilt},
\[
s_i^*
 \;\le\;
 \bcoef\T m + e^{-\vartheta \staleage_i}\,\bcoef\T(\tilde\xcov_i - m)
 + \tfrac12\bcoef\T V(\staleage_i)\bcoef ,
\]
so for firms whose last-seen signal was above the sector mean
($\bcoef\T\tilde\xcov_i > \bcoef\T m$), the naive score exceeds the
correct one by at least
$(1 - e^{-\vartheta \staleage_i})\,\bcoef\T(\tilde\xcov_i - m)
- \tfrac12\bcoef\T V(\staleage_i)\bcoef$, a quantity that is positive
and increasing in $\staleage_i$ once
$\bcoef\T(\tilde\xcov_i - m)$ clears the variance term, and the survival
tilt of \cref{lem:ch08:tilt} strictly enlarges the gap. Symmetrically,
below-mean stale firms are pushed up by mean reversion and down by the
tilt: their net bias is smaller and can take either sign.
\item \emph{Across the risk set}: $\staleage$ concentrates on
low-intensity firms (\cref{prop:ch08:internal}), so the distortion is
differential --- largest exactly where the naive score trusts old good
news. The characteristic failure is the \emph{stale champion}: a firm
that looked excellent at its last round years ago and has been silent
since is systematically over-ranked by LOCF plug-in, because its signal
should have been shrunk toward the prior and tilted down for the
silence, and was neither.
\end{enumerate}
\end{proposition}

\begin{proof}
(i) is arithmetic on \cref{eq:ch08:truelam}. (ii): the unconditional
(no-tilt) projection is \cref{eq:ch08:attenuation}, whose log is the
displayed bound; \cref{lem:ch08:tilt} applied to $h = \exp$ shows the
survival-conditioned expectation is strictly smaller. Subtracting from
$\tilde s_i = \bcoef\T\tilde\xcov_i$ gives the stated gap. (iii) combines
(ii) with \cref{prop:ch08:internal}(iii).
\end{proof}

\begin{example}[The size of the inversion]\label{ex:ch08:inversion}
Scalar covariate, $m = 0$, $\sigma_X = 1$, $\bcoef = 1$, covariate
half-life two years ($\vartheta = \log 2 / 2 \approx 0.347$ per year).
Firm $A$ last raised four years ago looking excellent, $\tilde x_A =
+2$; firm $B$ raised a quarter ago looking merely good, $\tilde x_B =
+1$. Naive LOCF plug-in scores $A$ far above $B$ ($2$ vs $1$). The
corrected log-scores of \cref{eq:ch08:attenuation}, before any survival
tilt: for $A$, $e^{-1.386}\cdot 2 + \tfrac12(1 - e^{-2.77}) \approx 0.50
+ 0.47 = 0.97$; for $B$, $e^{-0.087}\cdot 1 + \tfrac12(1 - e^{-0.173})
\approx 0.92 + 0.08 = 1.00$. The correction alone \emph{reverses} the
ranking, and the survival tilt of \cref{lem:ch08:tilt} --- four years of
silence versus one quarter --- pushes $A$ down further. At the risk-set
scale this is not an edge case: with inter-deal gaps of one to four
years and covariate half-lives of the same order, the attenuation factor
$e^{-\vartheta\staleage}$ ranges over $(0.25, 1)$ \emph{within a single
risk set}, so a model fitting one common slope is wrong for most of its
members at once.
\end{example}

\begin{warning}[Look-ahead through deal-time covariates is fatal]
\label{warn:ch08:lookahead}
The covariates in deal row $m$ are measured \emph{at or after} the deal:
revenue as reported in the round's diligence, raised-to-date
\emph{including} the round, valuation --- the outcome itself. Using row
$m$'s covariates as features for predicting deal $m$, or any as-of join
keyed on reference date rather than availability date, imports the label
into the features. This is the internal-covariate ``price sticker'' in
its most concrete form: the covariate value exists \emph{because} the
event happened. The failure mode is silent --- metrics inflate in
backtest and evaporate in deployment --- and no downstream remedy can
undo it. Policy: features at $t$ use deals strictly before $t$
(Protocol~\ref{prot:ch02:pit}); every loader ships a look-ahead audit that
verifies no feature path changes value at the timestamp of the label
event it feeds (\cref{ch08:sec:declaration}, code hooks). The synthetic
loader already enforces the strict inequality; the real-data loader must
inherit it verbatim.
\end{warning}

\section{Remedy R1: staleness-aware observed-data modeling}
\label{ch08:sec:r1}

The first remedy is not an attempt to undo the defect but to model what
we actually observe. Condition on the pair $(\tilde\xcov, \staleage)$
and fit
\begin{equation}\label{eq:ch08:r1}
\lam_i(t)
 \;=\;
 \exp\!\Big(
 \bcoef\T \tilde\xcov_i(t)
 \;+\; g\big(\staleage_i(t)\big)
 \;+\; \tilde\xcov_i(t)\T\, \Gamma_{\!g}\, \phi\big(\staleage_i(t)\big)
 \;+\; \delta\T \zfeat_i(t)
 \Big),
\end{equation}
with $g$ flexible (a spline or step function on an age grid with knots
at the empirical staleness quantiles), $\phi$ a low-dimensional basis in
the age whose interaction with $\tilde\xcov$ lets the covariate
\emph{slope} vary with staleness, and $\zfeat_i$ the other engineered
history covariates of \cref{def:ch02:covariates}(iii). By
\cref{prop:ch08:internal}(i) every term is $\mathcal{G}$-predictable, so
\cref{eq:ch08:r1} is a well-specified intensity model \emph{for the
observed filtration} and the likelihood \cref{eq:ch02:jacod} applies
with its full martingale estimation theory \citep{abgk1993,
andersengill1982}. The deep point is that this is not a heuristic patch:

\begin{proposition}[Projection: the observed-data intensity is a genuine
intensity]\label{prop:ch08:projection}
Let $N$ be a counting process adapted to a subfiltration
$\{\mathcal{G}_t\} \subseteq \{\Ft\}$ and let $\lam$ be its
$\F$-intensity. Then the $\mathcal{G}$-predictable projection
\[
\lam^{\mathcal{G}}(t) \;=\; \E\big[\lam(t) \st \mathcal{G}_{t^-}\big]
\]
is an intensity of $N$ with respect to $\{\mathcal{G}_t\}$: $N(t) -
\int_0^t \lam^{\mathcal{G}}(u)\dd u$ is a $\mathcal{G}$-local
martingale. Consequently the best-fitting element of a rich class of
functions of $(\tilde\xcov, \staleage, \Hist_{t^-})$ estimates the
projection of the true structural intensity onto the available
information --- fitting \cref{eq:ch08:r1} is estimating a well-defined
object, not compensating an error \citep[Ch.~14]{daley2008}
\citep[II.4]{abgk1993}.
\end{proposition}

\begin{proof}
For any bounded nonnegative $\mathcal{G}$-predictable $C$ (which is also
$\F$-predictable), the $\F$-intensity property gives
$\E\int_0^\infty C\dd N = \E\int_0^\infty C(u)\lam(u)\dd u$. Tower on
$\mathcal{G}_{u^-}$: since $C(u)$ is $\mathcal{G}_{u^-}$-measurable,
$\E[C(u)\lam(u)] = \E[C(u)\,\E(\lam(u)\st\mathcal{G}_{u^-})] =
\E[C(u)\lam^{\mathcal{G}}(u)]$, whence $\E\int C\dd N =
\E\int C\lam^{\mathcal{G}}$. This identity for all such $C$ is the
defining property of a $\mathcal{G}$-intensity (martingale
characterization via the monotone class theorem; the technical
identification of $\E[\,\cdot \st \mathcal{G}_{t^-}]$ with the
predictable projection is \citealp[Ch.~14]{daley2008}).
\end{proof}

\begin{corollary}[R1 is exactly specified in the OU world]
\label{cor:ch08:r1exact}
Under Assumption~\ref{asm:ch08:latent} with stationary drift and the
monotone-summary approximation of \cref{lem:ch08:tilt}, the projected
intensity is
$\exp\{(e^{-\Theta a}\bcoef)\T\tilde x + g_0(a) - d(\bcoef\T\tilde x,
a)\}$ with $g_0$ as in \cref{eq:ch08:attenuation} and $d \ge 0$ the
survival-tilt correction, a smooth function of the scalar signal and the
age. The class \cref{eq:ch08:r1} contains the first two terms exactly
(the interaction basis reproduces $e^{-\vartheta_k a}\bcoef_k$) and
approximates the third within the same $(\tilde\xcov\T\gamma,
\staleage)$ interaction structure. R1 is therefore not merely valid but
--- in the working model --- essentially well specified.
\end{corollary}

Two practical points before the caveats. First, the repository already
implements degenerate cases of \cref{eq:ch08:r1}: the risk-set ranker of
\modrepo{sector\_ranker.py} carries a months-since-last-funding
(cooldown) covariate --- a parametric $g$ inside the choice model --- and
the synthetic loader emits it point-in-time
(\code{log1p\_months\_since\_last\_funding} in \code{CHOICE\_FEATURES}).
D8.1 is the commitment that this stops being a per-model accident: the
pair $(\tilde\xcov, \staleage)$, with interactions where the model class
allows them, becomes mandatory plumbing in every stack. Second, an
identifiability warning that doubles as an interpretation warning: the
fitted $g$ absorbs \emph{everything} that varies with the gap ---
drift-variance inflation (\cref{prop:ch08:berkson}), mean-reversion
offsets (\cref{prop:ch08:attenuation}), the survival tilt, \emph{and}
genuine duration dependence of the hazard itself (post-round cooldown,
runway mechanics, the renewal structure of \cref{ch:hawkes}). These are
observationally confounded in $g$ by construction: the projection does
not distinguish ``covariates have drifted'' from ``the hazard truly
depends on time since last round.'' For prediction under the observed
filtration this is harmless --- the projection is the target
(\cref{prop:ch08:projection}) --- but it means $g$ has no structural
reading either, and separating the channels requires the state-space
machinery of \cref{ch08:sec:r3}, which supplies the drift-variance
component of $g$ from fitted dynamics and leaves the remainder as an
honest estimate of duration dependence.

What R1 gives up is exactly what the projection destroys. The fitted
$\bcoef$ in \cref{eq:ch08:r1} is a property of the \emph{observation
scheme composed with} the structural model: it mixes the true covariate
effect with mean reversion, drift variance, and the informativeness
tilt. It must not be read as ``the effect of revenue on funding
intensity,'' used for counterfactuals (``if this firm's revenue were
$20\%$ higher\dots'' --- the projection does not move that way), or
extrapolated to a vendor with a different update policy. And because
$(\tilde\xcov,\staleage)$ is internal, multi-horizon functionals require
joint simulation of events and the LOCF update process, never plug-in
survival formulas (\cref{ch08:sec:formal}). Within those limits, R1 is
committed, cheap, and always on --- decision D8.1.

\section{Remedy R2: supervised drift modeling and regression
calibration}\label{ch08:sec:r2}

R1 conditions on staleness; R2 tries to \emph{reverse} it. The key
observation is that the dataset labels its own drift: at every deal the
firm's true drifting covariates are re-observed. Every consecutive pair
of own deals $(\tau_m, \tau_{m'})$ of the same firm therefore yields a
supervised example
\[
\underbrace{\big(\xcov_i(\tau_m),\; a = \tau_{m'} - \tau_m,\;
\{\xcov^{\mathrm{mac}}_{\mathrm{pub}}(u)\}_{u \in (\tau_m, \tau_{m'}]},\;
\text{sector activity}\big)}_{\text{inputs}}
\;\longmapsto\;
\underbrace{\xcov_i(\tau_{m'})}_{\text{target}},
\]
a \emph{drift dataset} on which a conditional model $\hat m(\tilde x, a,
\cdot)$ with predictive variance $\hat V(\tilde x, a, \cdot)$ can be fit
--- linear-Gaussian (which is exactly the OU regression of
\cref{eq:ch08:oucond}, making R2 and R3 two estimators of one object) or
gradient-boosted with a variance head. The input design follows
\cref{eq:ch08:ou}: the macro path enters through summaries of the anchor
movement over the gap --- cumulative rate change, CPI and unemployment
changes between the two deal dates, as-of joined on publish dates per
Protocol~\ref{prot:ch02:pit}(a) --- and sector activity through trailing
deal counts, so that the linear-Gaussian fit is literally a discretized
regression estimate of $(\Theta, \Gamma, m_s, \Sigma_W)$, one
consecutive-pair likelihood term per row. The gradient-boosted variant
keeps the same inputs and relaxes the functional form; it is admitted
only if it beats the linear-Gaussian model on \emph{temporally} held-out
pairs (pairs whose second deal falls after the training era, per
\cref{prot:ch02:pit}(e)) on both mean error and variance calibration ---
a drift model whose $\hat V$ is uncalibrated poisons
\cref{eq:ch08:calib} through the exponent. Imputation then replaces the
stale value by $\hat\xcov_i(t) = \hat m(\tilde\xcov_i(t),
\staleage_i(t), \text{macro path})$, and the intensity uses regression
calibration:

\begin{proposition}[Plug-in of the conditional mean, with variance
correction]\label{prop:ch08:calibration}
Suppose $\xcov(t) \st \mathcal{G}_{t^-} \sim \mathcal{N}(\hat\xcov(t),
V(t))$. Then the projected intensity of \cref{prop:ch08:projection} is
\emph{exactly}
\begin{equation}\label{eq:ch08:calib}
\lam^{\mathcal{G}}(t)
 \;=\;
 \exp\!\Big(\bcoef\T\hat\xcov(t) + \tfrac12\,\bcoef\T V(t)\,\bcoef\Big).
\end{equation}
For a general conditional law with mean $\hat\xcov$ and variance $V$,
$\log\lam^{\mathcal{G}} = \bcoef\T\hat\xcov + \tfrac12\bcoef\T V \bcoef
+ O(\norm{\bcoef}^3)$ by the cumulant expansion: plugging in the
conditional \emph{mean} makes the first-order term exact, and the
variance factor $\exp(\tfrac12\bcoef\T V\bcoef)$ absorbs the second
order. Plugging in the raw $\tilde\xcov$, or a mode or median, leaves a
first-order error. This is regression calibration
\citep[Ch.~4]{carroll2006}, and \cref{eq:ch08:calib} is the log-linear
analogue of the induced-hazard formula of \citet{prentice1982} ---
including his caveat, which here is \cref{lem:ch08:tilt}: the
conditioning $\mathcal{G}_{t^-}$ includes the survival information, so
the honest $\hat\xcov$ is the \emph{tilted} mean, not the free-drift
mean.
\end{proposition}

\begin{proof}
The Gaussian case is the mgf again. In general, $\log\E[e^{\bcoef\T
\xcov}] = \sum_{r\ge1} c_r(\bcoef)/r!$ with first and second
cumulants $\bcoef\T\hat\xcov$ and $\bcoef\T V\bcoef$, the remainder
collecting third and higher cumulants of order $\norm{\bcoef}^3$.
\end{proof}

\begin{proposition}[Selection: the drift dataset is drift conditional on
success]\label{prop:ch08:selection}
A training pair exists only if a second deal occurred. Under the
monotone-summary approximation of \cref{lem:ch08:tilt} with
$\Lambda(s)=e^s$, the drift law represented in pairs with gap $a$ is
tilted relative to the free law by
$e^{s}\,e^{-a e^{s}}$ (event density at the endpoint), while the
population the imputation serves --- firms still waiting at age $a$ ---
is tilted by $e^{-a e^{s}}$ alone. The ratio of the two tilts is
$e^{s}$, increasing in $s = \bcoef\T\xcov(\tau + a)$, so the
pair-conditioned law dominates the survivor-conditioned law in the
likelihood-ratio order:
\[
\E\big[\bcoef\T\xcov(\tau+a) \st \text{pair with gap } a,\ \tilde x\big]
\;>\;
\E\big[\bcoef\T\xcov(\tau+a) \st \text{no event by } \tau + a,\ \tilde
x\big].
\]
A drift model fit naively on pairs therefore learns \emph{drift
conditional on funding success} and systematically over-predicts the
$\bcoef$-direction state of the silent firms it is deployed on ---
optimistic imputation, which re-creates the stale-champion over-ranking
of \cref{prop:ch08:ranking} through the back door.
\end{proposition}

\begin{proof}
The next-event density at gap $a$ given the endpoint summary is
$e^{s} e^{-a e^{s}}$ (hazard times survival), the no-event probability is
$e^{-a e^{s}}$; both statements are the conditional-Poisson formulas of
\cref{prop:ch08:internal}(iii). Their ratio $e^s$ is increasing, giving
likelihood-ratio dominance of the pair law over the survivor law, and
dominance implies the strict mean inequality for the nondegenerate
increasing functional $s$.
\end{proof}

Three remedies for the selection, in the order they will be tried.
\emph{Censoring-aware drift estimation}: treat pair generation as a
survival problem --- every firm silent since $\tau_m$ contributes a
right-censored drift spell $(\xcov_i(\tau_m), \staleage_i(\Tend),
\text{no endpoint})$, and the drift model is fit by a likelihood that
includes the survival factor, i.e.\ a small joint model of drift and
event (this is R4's machinery in miniature, and the honest default).
\emph{Inverse-probability weighting}: weight each pair by the inverse of
the estimated probability that a firm with those inputs produces a
second deal within the observed gap, taken from the discrete-hazard
bucket model of \cref{ch:buckets}; valid under the bucket model's own
correctness, cheap, and diagnosable. \emph{Exogenous restriction}: fit
drift only for covariates whose evolution is plausibly exogenous to the
funding process on the gap timescale --- headcount and revenue drift
with the business; raised-to-date, deal counts and last-round
descriptors are \emph{mechanical}, move only at events, and must never
be imputed (LOCF is exact for them, Assumption~\ref{asm:ch08:locf}). The
mechanical-vs-drifting classification of every PitchBook column is
maintained in \cref{app:columns} and is decision D8.3;
\cref{tab:ch08:columns} fixes the scheme on the representative columns,
including the third class that \cref{warn:ch08:lookahead} makes
necessary: columns that are \emph{outcomes} of the deal they ride on.

\begin{table}[htbp]
\centering
\small
\begin{tabular}{p{0.34\textwidth} p{0.13\textwidth} p{0.40\textwidth}}
\toprule
\textbf{Representative columns} & \textbf{Class} & \textbf{Staleness
treatment} \\
\midrule
\code{total\_raised\_to\_date}; deal count to date;
\code{last\_financing\_date}/\code{\_size}/\code{\_deal\_type};
\code{vc\_round} progression; firm age &
mechanical & LOCF exact by construction; recomputed point-in-time from
the event log; imputation forbidden (D8.3) \\
\addlinespace
\code{revenue}, \code{ebitda}, \code{gross\_profit},
\code{net\_income}, \code{number\_of\_employees}, growth since last
deal; valuation between rounds &
drifting & carried with $\staleage$ always (D8.1); eligible for R2
imputation with variance; modeled by the R3 state space \\
\addlinespace
\code{deal\_size}, \code{pre\_money\_valuation},
\code{post\_valuation}, up/down/flat flags of the \emph{current} round &
outcome-at-event & usable as covariates only strictly after their deal;
as predictors of their own deal they are look-ahead
(Warning~\ref{warn:ch08:lookahead}); as targets they belong to T4 \\
\addlinespace
\code{business\_status} &
drifting, laggy & exit-proxy territory of (C3); point-in-time only,
never applied retroactively \\
\bottomrule
\end{tabular}
\caption{The column classification of decision D8.3 on representative
PitchBook columns; the full dictionary lives in \cref{app:columns}.
``Mechanical'' columns are deterministic functionals of the observed
event history, for which LOCF is exact; ``drifting'' columns are the
carriers of defect (C5).}
\label{tab:ch08:columns}
\end{table}

\section{Remedy R3: an explicit covariate state space}
\label{ch08:sec:r3}

R2's linear-Gaussian special case deserves its own treatment, because it
is simultaneously the honest uncertainty model and the supplier of the
$\Sigma(a)$, $V(a)$ objects that R1 and R2 consume. Model the drifting
components per firm by \cref{eq:ch08:ou} --- OU toward a sector mean
with macro loadings --- observed at own deals with reporting noise:
$y_{i,m} = \xcov_i(\tau_m) + \eta_m$, $\eta_m \sim \mathcal{N}(0, R)$
(PitchBook financials are reported figures, not audited truth; $R = 0$
recovers exact observation).

\begin{lemma}[OU prediction and deal-time update]\label{lem:ch08:ou}
Under \cref{eq:ch08:ou} with constant mean $m$ (time-varying $m_i(u)$
adds the deterministic integral $\int_\tau^t e^{-\Theta(t-u)}\Theta\,
m_i(u)\dd u$, computable from the published macro path), the exact
recursions for the conditional law $\xcov_i(t) \st \mathcal{G}_t \sim
\mathcal{N}(\hat\xcov_i(t), P_i(t))$ are:
\begin{align}
\text{predict over } a = t - \tau:\quad
&\hat\xcov_i(t) \;=\; m + e^{-\Theta a}\big(\hat\xcov_i(\tau) - m\big),
\nonumber\\
&P_i(t) \;=\; e^{-\Theta a} P_i(\tau)\, e^{-\Theta a}
 + V(a), \qquad
 V(a)_{kk} = \sigma^2_{X,k}\big(1 - e^{-2\vartheta_k a}\big),
\label{eq:ch08:kalmanpredict}\\[2pt]
\text{update at an own deal with report } y:\quad
&K = P\,(P + R)^{-1}, \qquad
\hat\xcov^{+} = \hat\xcov + K\,(y - \hat\xcov), \qquad
P^{+} = (I - K)\,P.
\label{eq:ch08:kalmanupdate}
\end{align}
With $R = 0$ the update collapses to $\hat\xcov^+ = y$, $P^+ = 0$, and
the filter between deals is pure prediction: $\hat\xcov_i(t)$ is exactly
the shrunk mean of \cref{eq:ch08:oucond} and $P_i(t) = V(\staleage_i(t))$
is the explicit $\Sigma(a)$ that
\cref{prop:ch08:berkson,prop:ch08:attenuation} postulated.
\end{lemma}

\begin{proof}
The predict step is \cref{eq:ch08:oucond} applied to the current
posterior (linear-Gaussian closure under the OU transition); the update
step is the standard Kalman gain for a linear observation with Gaussian
noise \citep{kalman1960,harvey1989}. The $R=0$ limit is immediate.
\end{proof}

The state-space parameters $(\Theta, \Sigma_W, m_s, \Gamma, R)$ are
estimable by exact Gaussian likelihood on the drift dataset of
\cref{ch08:sec:r2} --- with the same selection caveat, treated by the
same censoring-aware likelihood: \cref{prop:ch08:selection} applies to
$\hat\Theta$ too, and ignoring it biases mean-reversion estimates toward
the reversion experienced by successful firms. Estimation must pool: the
median tracked firm contributes a handful of deals, far too few to
identify per-firm dynamics, so $(\Theta, \Sigma_W)$ are shared at sector
level and $\Gamma$ globally, fitted by summing pair likelihoods across
firms; the drifting covariates enter on log scales (log revenue, log
headcount), where the Gaussian working model is tenable and where the
loaders already live, with missing and undisclosed reports handled by
the rules of \cref{app:columns}, never ad hoc. Firm-level heterogeneity
beyond the pooled parameters is precisely the shared random effect that
R4 would add --- another sense in which the ladder is nested rather than
branching. What R3 buys beyond R2:
(i) $V(a)$ in closed form, feeding R1's offset $g_0$ and R2's variance
correction \cref{eq:ch08:calib} instead of leaving them as free spline
coefficients; (ii) calibrated per-firm uncertainty that grows honestly
to the stationary variance as $\staleage \to \infty$ --- the model
\emph{knows} it knows nothing about a firm silent for a decade, which is
precisely the epistemic behavior the ranking layer needs; (iii) a
principled bridge to the doubly-stochastic intensity tradition ---
default intensities driven by stochastic covariates with fitted dynamics
--- whose viability at scale is established in credit risk
\citep{duffie2007}. The cost is the linear-Gaussian straitjacket, which
is why R3 enters the ladder as the variance model \emph{inside} R2
rather than as a replacement for the gradient-boosted drift mean.

\section{Remedy R4: joint modeling, gated by a measured trigger}
\label{ch08:sec:r4}

The statistically complete treatment is a joint model of the
longitudinal covariates and the event intensity with shared random
effects \citep{rizopoulos2012,tsiatisdavidian2004}: latent per-firm
effects $b_i$ drive both the covariate trajectory ($\xcov_i(t) =
\mu(t) + \Phi(t) b_i + \epsilon_i(t)$) and the intensity
($\lam_i(t) = \exp(\bcoef\T\xcov_i(t) + \alpha\T b_i + \cdots)$), and the
likelihood integrates $b_i$ out over both data streams jointly. This is
the construction that makes the informative observation scheme part of
the likelihood rather than a correction: the two-stage shortcut --- fit
the longitudinal model, plug fitted trajectories into the intensity ---
was shown biased precisely when the event process informs the covariate
observations \citep{tsiatis1995}, which is our
\cref{prop:ch08:internal}(iii) verbatim, and the bias direction is the
one \cref{prop:ch08:selection} derives: trajectories fitted to
success-selected observations are optimistic for the silent. The joint
likelihood repairs this by making the no-event stretches pay their
survival factor inside the same integral that fits the trajectory ---
R2's censoring-aware drift estimation and R3's censoring-aware
state-space likelihood are exactly the first two terms of this
construction, which is why the ladder converges to R4 rather than being
overturned by it. What R4 adds beyond them is the shared random effect:
persistent firm-level frailty entering both streams, so that a firm's
covariate \emph{level and trajectory shape} inform its intensity beyond
the current imputed value \citep{tsiatisdavidian2004,rizopoulos2012}.
The price is steep: numerical integration over random effects for every
firm ($\sim$7{,}000 tracked firms $\times$ quadrature nodes $\times$
optimizer iterations), delicate identifiability between $\alpha$,
$\bcoef$ and the drift parameters --- the frailty-vs-observation
confounding of \cref{app:results}'s branching-inflation verdict returns
here in longitudinal form --- and a model class fragile enough that its
failure modes are harder to audit than the defect it fixes. R4 is
therefore \emph{escalation, not default}, and it is gated:

\begin{protocol}[Attenuation-gap test, the R4 trigger]
\label{prot:ch08:gaptest}
On the synthetic rehearsal world with drifting firm covariates, LOCF
observation, and known structural $\bcoef$: (i) run the committed
pipeline R1 $+$ R2 (with R3's variance); (ii) measure the attenuation
ratio $\hat r$ = slope of the fitted log-intensity scores regressed on
the oracle scores $\bcoef\T\xcov_i(t)$ over the risk set, and the top-$K$
ranking gap between fitted and oracle covariates; (iii) the trigger
fires if the $95\%$ interval for $\hat r$ excludes $[0.9, 1.1]$ \emph{and}
the oracle-covariate ablation attributes a material top-$K$ loss to
residual staleness (not to (C3), which is measured separately per
\cref{ch:data}), in two consecutive evaluation campaigns. Only a fired
trigger authorizes the R4 build; the same harness then has to show R4
closing the measured gap on rehearsal data before any real-data use.
\end{protocol}

\section{The bucket-model face, and what the rehearsal measured}
\label{ch08:sec:buckets}

\Cref{ch:buckets} attacks the horizon target T3 with gradient boosting
on a bucketed panel whose feature rows read, in words: ``it has been
$\staleage$ since this firm's last deal, and these are its last-known
covariates.'' That is $(\tilde\xcov, \staleage)$ --- and the trees'
recursive partitioning of the joint feature space fits arbitrary
piecewise-constant surfaces in $(\tilde\xcov, \staleage)$, i.e.\
nonparametric estimates of exactly the projected intensity that
\cref{cor:ch08:r1exact} derived: offset $g_0(\staleage)$,
age-attenuated slopes $e^{-\Theta\staleage}\bcoef$, and the survival-tilt
correction, none of which need to be prespecified
\citep{friedman2001,chenguestrin2016}. Gradient boosting on
$(\tilde\xcov, \staleage)$ \emph{is} R1, implemented nonparametrically.
This is the principled reason the discrete-time stack tolerates
staleness better than a rigid log-linear form fitted without age terms:
the rigid form suffers the pooled attenuation and signed
omitted-variable bias of \cref{cor:ch08:pooled}, while the trees simply
learn the interaction surface. It is also a plausible component of the
fired B1 test --- the parity XGBoost beating the linear conditional
logit by $+0.22$ nats and $+9$pp top-5 on regime A
(\cref{app:results}) --- though the rehearsal has not yet isolated how
much of that gap is staleness interaction versus other nonlinearity;
the isolation run is scheduled below.

The rehearsal evidence on the attenuation mechanics themselves comes
from the covariate-noise campaign exp22, catalogued in
\cref{app:results} and inherited here without re-derivation: relative
Gaussian noise on the covariates at fit and prediction time --- the
\emph{classical-error benchmark} of \cref{ch08:sec:attenuation}, not yet
informative staleness --- attenuated the fitted choice weights by
$\times 0.88$ at $20\%$ noise ($\times 0.93$ at $15\%$) and eroded the
stage-1 $\bcoef$ correlation from $0.67$ to $0.53$, while prediction was
nearly untouched (stage-1 covariate lift $\approx 1.0$ nats throughout;
top-5 moved $0.446 \to 0.439$). Two lessons transfer. First, the
attenuation arithmetic of \cref{prop:ch08:attenuation} is not
hypothetical --- it fires at the predicted order of magnitude on data of
our shape, and it hits \emph{interpretation} (structural readings of
$\bcoef$, cross-regime transfer) long before it hits ranking or
likelihood. Second, exp22's noise was exogenous and uninformative ---
independent of the risk set --- so it measures the benchmark, not the
twist: the differential, informative distortion of
\cref{prop:ch08:ranking} is precisely what exp22 could \emph{not} see.
The planned follow-up (same harness, drifting covariates with LOCF ages
emitted by the loader, ranking bias reported by staleness decile) closes
that gap and doubles as the standing harness for
Protocol~\ref{prot:ch08:gaptest}.

Because aggregate metrics demonstrably dilute exactly the damage this
chapter is about --- exp22's headline prediction numbers barely moved
while the fitted structure decayed --- the chapter installs a standing
reporting discipline:

\begin{protocol}[Staleness accounting in every evaluation]
\label{prot:ch08:account}
Every model evaluation report in the repository includes: (a) the
distribution of $\staleage$ over the evaluation risk sets (quantiles,
and the share exceeding one and two estimated covariate half-lives once
R3 supplies $\hat\Theta$); (b) the headline metrics stratified by
staleness quartile of the scored firm; (c) for ranking deliverables, the
share of top-$K$ selections whose drifting covariates are older than one
year, reported next to the hit rate. Rationale:
\cref{prop:ch08:ranking} localizes the damage in the stale strata, so a
model can win on average while being systematically wrong about the
firms it is most confidently promoting from old information --- and the
consuming action (a diligence short list) is exposed to exactly those
firms. The stratified report is the cheap standing detector of residual
(C5) damage and the attribution input to
Protocol~\ref{prot:ch08:gaptest}(iii).
\end{protocol}

\section{Defect declaration and commitments}\label{ch08:sec:declaration}

Per decision D2.1, the models of this chapter declare their position
against the ledger of \cref{tab:ch02:defects}. \textbf{(C5)} is
\emph{owned here}: R1 corrects the fitting bias exactly in the OU
working model (\cref{cor:ch08:r1exact}) and validly in general
(\cref{prop:ch08:projection}); R2/R3 additionally correct the ranking
bias under a correct, selection-aware drift model; residual
vulnerability is the survival-tilt approximation and drift-model error,
monitored by Protocol~\ref{prot:ch08:gaptest}. \textbf{(C1)}: robust ---
all constructions are point-in-time and censored spells enter the drift
likelihood as censored, per \cref{ch08:sec:r2}; dropping them would
reimport \cref{prop:ch08:selection} with the sign derived there.
\textbf{(C2)}: robust by construction --- $(\tilde\xcov, \staleage)$ is
defined only after first observation; pre-entry firms are the newcomer
problem of \cref{ch:universe}. \textbf{(C3)}: \emph{vulnerable, with a
signed interaction that must be stated}. An unrecorded-dead firm
accumulates maximal staleness with frozen covariates: (C3) keeps it in
the risk set, (C5) keeps its last good numbers on its face. The two
defects compound on the same firms --- the stale champion of
\cref{prop:ch08:ranking} is, in the worst case, dead --- and both
corrupt the risk set's information state in the same direction:
overstated scores for long-silent firms, hence jointly worse for
ranking than either alone. The staleness terms $g(\staleage)$ partially
absorb (C3) (large age predicts death as well as drift), which improves
ranking but permanently confounds ``quiet'' with ``gone'': no structural
reading of $g$ is permitted while (C3) is open. \textbf{(C4)}: robust
--- staleness ages are measured in weeks-to-years and deal-date jitter
is days; aggregation does not touch the LOCF map. \textbf{(C6)}:
out-of-universe firms are the $\staleage \to \infty$ limit --- no
covariates at all --- and are handled by the outside-option construction
of \cref{ch:universe}, which this chapter's $g$ should be checked
against for continuity (a firm silent for a decade should not score very
differently from an unseen firm).

\begin{decision}{D8.1 --- staleness age is a mandatory feature}
From this document forward, every firm-level model in every stack ---
event-time intensities, the risk-set ranker, the bucketed panel, mark
regressions --- includes the staleness age $\staleage_i(t)$ (and, where
the model class permits, its interactions with $\tilde\xcov_i$) among
its features. Rationale: $\staleage$ is free, predictable
(\cref{prop:ch08:internal}), and carries the offset and attenuation
structure of \cref{prop:ch08:attenuation}; omitting it produces the
signed biases of \cref{cor:ch08:pooled,prop:ch08:ranking}. Reversal
trigger: none --- a fitted model may drive its age terms to zero, but
the feature may not be withheld from it.
\end{decision}

\begin{decision}{D8.2 --- the remedy ladder}
Committed sequence: R1 (staleness-aware observed-data modeling,
\cref{eq:ch08:r1}) now and always on; R2 (drift dataset, supervised
drift model, regression calibration per \cref{prop:ch08:calibration},
selection handled censoring-aware first) next quarter; R3 (OU/Kalman
state layer, \cref{lem:ch08:ou}) folded in as R2's variance model rather
than a separate stack; R4 (joint model) only if
Protocol~\ref{prot:ch08:gaptest} fires. Rationale: the ladder is cost-ordered,
each rung consumes the previous rung's outputs, and the rehearsal
evidence says prediction is initially robust (exp22) --- so the
expensive rungs must prove their marginal value on the measured gap, not
on principle. Reversal trigger: a fired gap test promotes R4; a gap
test passing twice consecutively with R1 alone demotes R2/R3 to
research status.
\end{decision}

\begin{decision}{D8.3 --- mechanical vs.\ drifting column classification}
Every PitchBook covariate column is classified in the data dictionary
(\cref{app:columns}) as \emph{mechanical} (deterministic functional of
the observed event history: raised-to-date, deal counts, last-round
descriptors, ownership accumulations --- LOCF exact, imputation
forbidden) or \emph{drifting} (revenue, EBITDA, net income, employees,
growth rates, latent valuation --- LOCF stale, carried with
$\staleage$, eligible for R2/R3 treatment). The classification is
maintained under vendor schema changes and consumed by the loaders and
the drift-dataset builder. Rationale: Assumption~\ref{asm:ch08:locf} shows (C5)
binds only on drifting columns, and imputing a mechanical column is not
an approximation but an error. Reversal trigger: none; reclassification
of individual columns requires a data-dictionary change with review.
\end{decision}

\begin{codehooks}
Planned module \modrepo{covstate/}: \code{pairs.py} --- the drift-dataset
builder (consecutive own-deal pairs with gap, macro path and
sector-activity summaries; right-censored spells emitted with flags per
\cref{ch08:sec:r2}; selection diagnostics from
\cref{prop:ch08:selection}); \code{ou.py} --- OU transition and
$V(\staleage)$ in closed form, Kalman predict/update of
\cref{lem:ch08:ou}, censoring-aware Gaussian likelihood for $(\Theta,
\Sigma_W, m_s, \Gamma, R)$; \code{impute.py} --- the imputation API
returning $(\hat\xcov_i(t), V_i(t))$ and the calibrated intensity factor
$\exp(\tfrac12\bcoef\T V\bcoef)$ of \cref{eq:ch08:calib}. Loaders:
\code{synthetic/loaders.py} already enforces strictly-prior LOCF and
carries a months-since-last-funding feature; it and the real-data
adapter \modrepo{data/pitchbook.py} must emit the pair $(\tilde\xcov_i,
\staleage_i)$ explicitly for every stack (D8.1), keyed to the
mechanical/drifting classification of \cref{app:columns} (D8.3).
Audits (\modrepo{data/audits.py}): the look-ahead shift test of
\cref{warn:ch08:lookahead} --- no feature path may change value at the
timestamp of the label event it feeds --- plus a staleness-age
consistency check ($\staleage$ recomputed from the event list must match
the emitted feature). The gap-test harness of Protocol~\ref{prot:ch08:gaptest}
extends the exp22 campaign driver with LOCF-age worlds and
staleness-decile ranking reports.
\end{codehooks}
